% !TeX root = main.tex
% !TeX spellcheck = pt_BR
% ===========================================================
% Testes de Hipóteses II: RVG e testes UMP
% Fonte: 10 - Testes de Hipóteses 2.pdf
% ===========================================================

% ---------- Introdução ----------
Até o momento, estudamos testes envolvendo hipóteses simples, o que não é usual em situações reais. Agora vamos dar ênfase a testes que envolvem hipóteses compostas.

Vamos trabalhar sob a suposição de que temos uma amostra aleatória oriunda de uma população com f.d.p. (ou f.p., embora bem menos)
\begin{equation}
f(\v{x}\mid\theta), \quad \theta\in\Theta.
\end{equation}
Nosso interesse será testar hipóteses $\mathcal{H}_0:\theta\in\Theta_0$ \emph{vs.} $\mathcal{H}_1:\theta\in\Theta_1$, em que $\Theta_0\subset\Theta$, $\Theta_1\subset\Theta$, $\Theta_0\cup\Theta_1=\Theta$ e $\Theta_0\cap\Theta_1=\emptyset$.
% Correspondência: 10 - Testes de Hipóteses 2.pdf, p. 3/95

% ---------- Definição 1 (RVG) ----------
\begin{definition}[Razão de verossimilhanças generalizada -- RVG]
Seja $L(\theta\mid\v{x})$ a função de verossimilhança para uma amostra $X_1,X_2,\ldots,X_n$, que tem f.d.p. (ou f.p.) conjunta $f(\v{x}\mid\theta)$, $\theta\in\Theta$. Chamamos de \uline{razão de verossimilhanças generalizada} a razão
\begin{equation}
\lambda=\lambda(\v{x})=\frac{\sup_{\theta\in\Theta_0}L(\theta\mid\v{x})}{\sup_{\theta\in\Theta}L(\theta\mid\v{x})}.
\end{equation}
\end{definition}
% Correspondência: 10 - Testes de Hipóteses 2.pdf, p. 5/95

% ---------- Texto: algumas ponderações sobre a RVG ----------
Observe que $\lambda$ é uma função da amostra observada $\v{x}$. Quando a RVG envolver as v.a.'s $\v{X}$, então a denotaremos por $\Lambda=\lambda(\v{X})$. Por vezes, podemos denotar $\Lambda_n$ ou $\lambda_n$, quando quisermos enfatizar o tamanho amostral.

Note também que $\Lambda$ é uma estatística, por ser uma função de $\v{X}$ que não depende de quantidades desconhecidas.

A partir da Definição 1, temos que $0\leq\lambda\leq1$. Uma vez que numerador e denominador não podem ser negativos, $\lambda\geq0$. Por outro lado, como $\Theta_0\subset\Theta$, então $\sup_{\theta\in\Theta_0}L(\theta\mid\v{x})\leq\sup_{\theta\in\Theta}L(\theta\mid\v{x})$, que implica $\lambda\leq1$.

Valores pequenos de $\lambda$ indicam evidências contra $\mathcal{H}_0$.
% Correspondência: 10 - Testes de Hipóteses 2.pdf, p. 6/95

Uma generalização natural do teste da razão de verossimilhanças, visto para testar hipóteses simples, seria
\begin{equation}
\lambda^*(\v{x})=\frac{\sup_{\theta\in\Theta_0}L(\theta\mid\v{x})}{\sup_{\theta\in\Theta_1}L(\theta\mid\v{x})}.
\end{equation}
Por ser mais simples de se calcular, o teste da razão de verossimilhanças generalizada é adotado na forma de $\lambda(\v{x})$, exibida na Definição 1.

Vale destacar que os testes $\lambda(\v{x})$ e $\lambda^*(\v{x})$ estão relacionados da forma
\begin{equation}
\lambda(\v{x})=\min\{1,\lambda^*(\v{x})\}.
\end{equation}
% Correspondência: 10 - Testes de Hipóteses 2.pdf, p. 7/95

% ---------- Observação 1 ----------
\begin{note}
	\leavevmode
	\begin{outline}[enumerate]
		\1 O parâmetro $\theta$ pode ser um escalar ou um vetor. Além disso, numerador e denominador de $\lambda(\v{x})$ correspondem aos EMV's de $\theta$ em $\Theta_0$ e $\Theta$, respectivamente, aplicados à função de verossimilhança.
	\end{outline}
\end{note}
% Correspondência: 10 - Testes de Hipóteses 2.pdf, p. 8/95

% ---------- Texto: princípio do TRVG e pontos negativos ----------
Valores $\lambda$ da estatística $\Lambda$ são usados para a tomada de decisão em um teste do tipo $\mathcal{H}_0:\theta\in\Theta_0$ \emph{vs.} $\mathcal{H}_1:\theta\in\Theta_1$. Para isso, utilizamos o princípio do teste da razão de verossimilhanças generalizada: rejeitar $\mathcal{H}_0$ se, e somente se, $\lambda\leq\lambda_0$, em que $\lambda_0$ é uma constante tal que $0<\lambda_0<1$. Essa constante é, geralmente, especificada fixando o tamanho do teste.

Alguns pontos negativos do TRVG:
\begin{outline}
\1 Às vezes é difícil determinar $\sup L(\theta\mid\v{x})$;
\1 Às vezes é difícil determinar a distribuição de $\Lambda$, fundamental para calcular o poder do teste.
\end{outline}
% Correspondência: 10 - Testes de Hipóteses 2.pdf, p. 9/95

% ---------- Observação 2 ----------
\begin{note}
	\leavevmode
	\begin{outline}[enumerate]
		\1 O desenvolvimento de testes da RVG geralmente começa de uma forma um tanto complicada, até chegar em resultados mais simples.
	\end{outline}
\end{note}
% Correspondência: 10 - Testes de Hipóteses 2.pdf, p. 11/95

% ---------- Observação 3 ----------
\begin{note}
	\leavevmode
	\begin{outline}[enumerate]
		\1 Por construção, testes da RVG dependem da amostra apenas através de estatísticas suficientes (mínimas).
	\end{outline}
\end{note}
% Correspondência: 10 - Testes de Hipóteses 2.pdf, p. 11/95

% ---------- Definição 2 (TUMP) ----------
\begin{definition}[Testes uniformemente mais poderosos -- TUMP]
Um teste $\Upsilon^*$ de $\mathcal{H}_0:\theta\in\Theta_0$ \emph{vs.} $\mathcal{H}_1:\theta\in\Theta_1$ é chamado de \uline{teste uniformemente mais poderoso} de tamanho $\alpha$ se, e somente se,
\begin{outline}
\1 (i) $\sup_{\theta\in\Theta_0}\pi_{\Upsilon^*}(\theta)=\alpha$;
\1 (ii) $\pi_{\Upsilon^*}(\theta)\geq\pi_\Upsilon(\theta)$, para todo $\theta\in\Theta_1$ e para qualquer teste $\Upsilon$ com tamanho menor do que ou igual a $\alpha$.
\end{outline}
\end{definition}
% Correspondência: 10 - Testes de Hipóteses 2.pdf, p. 13/95

% ---------- Observação 4 ----------
\begin{note}
	\leavevmode
	\begin{outline}[enumerate]
		\1 O advérbio \emph{uniformemente} está associado aos parâmetros $\theta\in\Theta_1$.
	\end{outline}
\end{note}
% Correspondência: 10 - Testes de Hipóteses 2.pdf, p. 13/95

% ---------- Texto de transição ----------
O Lema de Neyman-Pearson é fundamental para mostrarmos que um teste é UMP em casos como esse.

Há formas de estender o resultado apresentado para casos mais gerais. Os Teoremas 1 e 2 apresentados a seguir são fundamentais para isso.
% Correspondência: 10 - Testes de Hipóteses 2.pdf, p. 15/95

% ---------- Teorema 1 ----------
\begin{theorem}
Sejam $X_1,X_2,\ldots,X_n$ uma A.A. oriunda de uma população com f.d.p. $f(x\mid\theta)$, $\theta\in\Theta$, em que $\Theta$ é algum intervalo. Suponha que $f(x\mid\theta)$ pertence à família exponencial, ou seja,
\begin{equation}
f(x\mid\theta)=\exp\{\eta(\theta)T(x)-B(\theta)\}h(x).
\end{equation}
Seja $T^*(\v{x})\equiv\sum_{i=1}^n T(x_i)$.
\begin{outline}
\1 (i) Se $\eta(\theta)$ for uma função monótona crescente de $\theta$, e se existir uma constante $k^*$ tal que $\P_{\theta_0}(T^*(\v{X})>k^*)=\alpha$, então o teste $\Upsilon^*$ com região crítica $\mathcal{C}^*=\{\v{x}\in\mathbb{R}^n\mid T^*(\v{x})>k^*\}$ será um teste UMP de tamanho $\alpha$ para $\mathcal{H}_0:\theta\leq\theta_0$ \emph{vs.} $\mathcal{H}_1:\theta>\theta_0$ ou $\mathcal{H}_0:\theta=\theta_0$ \emph{vs.} $\mathcal{H}_1:\theta>\theta_0$;
\1 (ii) Se $\eta(\theta)$ for uma função monótona decrescente de $\theta$, e se existir uma constante $k^*$ tal que $\P_{\theta_0}(T^*(\v{X})<k^*)=\alpha$, então o teste $\Upsilon^*$ com região crítica $\mathcal{C}^*=\{\v{x}\in\mathbb{R}^n\mid T^*(\v{x})<k^*\}$ será um teste UMP de tamanho $\alpha$ para $\mathcal{H}_0:\theta\leq\theta_0$ \emph{vs.} $\mathcal{H}_1:\theta>\theta_0$ ou $\mathcal{H}_0:\theta=\theta_0$ \emph{vs.} $\mathcal{H}_1:\theta>\theta_0$.
\end{outline}
\end{theorem}
% Correspondência: 10 - Testes de Hipóteses 2.pdf, p. 16/95

% ---------- Definição 3 (Razão de verossimilhanças monótona) ----------
\begin{definition}[Razão de verossimilhanças monótona]
Uma família de f.d.p.'s $\{f(x\mid\theta)\mid\theta\in\Theta,\ \Theta\text{ intervalo}\}$ é dita ter \uline{razão de verossimilhanças monótona} se, de uma amostra $X_1,X_2,\ldots,X_n$, houver uma estatística $T\equiv T(\v{X})$ tal que $\dfrac{L(\theta'\mid\v{x})}{L(\theta''\mid\v{x})}$ seja uma função não crescente de $T(\v{x})$, para todo $\theta'<\theta''$, ou uma função não decrescente de $T(\v{x})$, para todo $\theta'<\theta''$.
\end{definition}
% Correspondência: 10 - Testes de Hipóteses 2.pdf, p. 18/95

% ---------- Observação 5 ----------
\begin{note}
	\leavevmode
	\begin{outline}[enumerate]
		\1 A razão de verossimilhanças monótona é apenas uma razão de duas funções de verossimilhança, e não uma razão de verossimilhanças generalizada.
	\end{outline}
\end{note}
% Correspondência: 10 - Testes de Hipóteses 2.pdf, p. 18/95

% ---------- Teorema 2 ----------
\begin{theorem}
Sejam $X_1,X_2,\ldots,X_n$ uma A.A. oriunda de uma população com f.d.p. $f(x\mid\theta)$, em que $\Theta$ é um intervalo. Suponha que a família de f.d.p.'s $\{f(x\mid\theta)\mid\theta\in\Theta\}$ tenha razão de verossimilhanças monótona na estatística $T(\v{X})$. Nesse caso:
\begin{outline}
\1 (i) Se a razão de verossimilhanças monótona for não decrescente em $T(\v{X})$, e se $k^*$ for tal que $\P_{\theta_0}(T(\v{X})<k^*)=\alpha$, então o teste correspondendo à região crítica $\mathcal{C}^*=\{\v{x}\in\mathbb{R}^n\mid T(\v{x})<k^*\}$ será UMP de tamanho $\alpha$ para $\mathcal{H}_0:\theta\leq\theta_0$ \emph{vs.} $\mathcal{H}_1:\theta>\theta_0$;
\1 (ii) Se a razão de verossimilhanças monótona for não crescente em $T(\v{x})$, e se $k^*$ for tal que $\P_{\theta_0}(T(\v{X})>k^*)=\alpha$, então o teste correspondendo à região crítica $\mathcal{C}^*=\{\v{x}\in\mathbb{R}^n\mid T(\v{x})>k^*\}$ será UMP de tamanho $\alpha$ para $\mathcal{H}_0:\theta\leq\theta_0$ \emph{vs.} $\mathcal{H}_1:\theta>\theta_0$.
\end{outline}
\end{theorem}
% Correspondência: 10 - Testes de Hipóteses 2.pdf, p. 21/95

% ---------- Observação 6 ----------
\begin{note}
	\leavevmode
	\begin{outline}[enumerate]
		\1 As hipóteses utilizadas nos Teoremas 1 e 2 são do tipo $\mathcal{H}_0:\theta\leq\theta_0$ \emph{vs.} $\mathcal{H}_1:\theta>\theta_0$. Contudo, esses mesmos teoremas podem ser reformulados para hipóteses na forma $\mathcal{H}_0:\theta\geq\theta_0$ \emph{vs.} $\mathcal{H}_1:\theta<\theta_0$, desde que os sinais de desigualdade das regiões críticas sejam invertidos.
	\end{outline}
\end{note}
% Correspondência: 10 - Testes de Hipóteses 2.pdf, p. 23/95

% ---------- Observação 7 ----------
\begin{note}
	\leavevmode
	\begin{outline}[enumerate]
		\1 Os referidos teoremas consideram apenas hipóteses unilaterais.
	\end{outline}
\end{note}
% Correspondência: 10 - Testes de Hipóteses 2.pdf, p. 23/95

% ---------- Texto: resumo da ópera ----------
Vimos que testes UMP existem para hipóteses unilaterais, quando a razão de verossimilhanças é monótona em alguma estatística.

Há vários problemas de testes de hipóteses em que não há um teste UMP.

Podemos pensar, portanto, em restringir a classe de testes, com o intuito de determinar um teste ótimo. Nesse caso, uma alternativa interessante corresponde aos testes não viesados.
% Correspondência: 10 - Testes de Hipóteses 2.pdf, p. 23/95

% ---------- Definição 4 (Teste não viesado) ----------
\begin{definition}
Um teste $\Upsilon$ da hipótese nula $\mathcal{H}_0:\theta\in\Theta_0$ \emph{vs.} a hipótese alternativa $\mathcal{H}_1:\theta\in\Theta_1$ é chamado de \uline{teste não viesado} se, e somente se,
\begin{equation}
\sup_{\theta\in\Theta_0}\pi_\Upsilon(\theta)\leq\inf_{\theta\in\Theta_1}\pi_\Upsilon(\theta).
\end{equation}
\end{definition}
% Correspondência: 10 - Testes de Hipóteses 2.pdf, p. 26/95

% ---------- Texto: testes não viesados ----------
Para testes não viesados, a probabilidade de rejeitar $\mathcal{H}_0$, quando essa hipótese for verdadeira, nunca será superior à probabilidade de rejeitá-la quando ela for falsa.

Se um teste não viesado for UMP, então ele será chamado de teste não viesado uniformemente mais poderoso.

Não entraremos em detalhe sobre essa classe de testes.
% Correspondência: 10 - Testes de Hipóteses 2.pdf, p. 26/95

% ---------- Contextualização: amostragem de uma população normal ----------
Suponhamos que a amostra aleatória $X_1,X_2,\ldots,X_n$ seja oriunda de uma população $X\sim\normdist{\mu,\sigma^2}$.

Vamos analisar testes de hipóteses para a média $\mu$ e para a variância $\sigma^2$.
% Correspondência: 10 - Testes de Hipóteses 2.pdf, p. 28/95

% ---------- Testes para a média: introdução ----------
O mais comum é estarmos interessados em testes unilaterais ou bilaterais.

No caso dos testes unilaterais, vamos abordar $\mathcal{H}_0:\mu\leq\mu_0$ \emph{vs.} $\mathcal{H}_1:\mu>\mu_0$.

Há dois casos a serem considerados:
\begin{outline}
\1 $\sigma^2$ conhecido;
\1 $\sigma^2$ desconhecido.
\end{outline}
% Correspondência: 10 - Testes de Hipóteses 2.pdf, p. 29/95

% ---------- Teorema/desenvolvimento: H0: mu <= mu0 vs H1: mu > mu0, sigma^2 conhecido ----------
\begin{theorem}[$\mathcal{H}_0:\mu\leq\mu_0$ \emph{vs.} $\mathcal{H}_1:\mu>\mu_0$ -- $\sigma^2$ conhecido]
Sejam $X_1,X_2,\ldots,X_n$ uma A.A. de uma população $X\sim\normdist{\mu,\sigma^2}$, com $\sigma^2$ conhecido. Nesse caso, o espaço paramétrico corresponderá ao conjunto dos reais. Além disso,
\begin{align}
f(x\mid\mu) &= \frac{1}{\sqrt{2\pi\sigma^2}}\exp\left\{-\frac{1}{2\sigma^2}(x-\mu)^2\right\}\notag\\
&= \frac{1}{\sqrt{2\pi\sigma^2}}\exp\left\{-\frac{1}{2\sigma^2}x^2+\frac{\mu}{\sigma^2}x-\frac{\mu^2}{2\sigma^2}\right\}\notag\\
&= \exp\left\{\frac{\mu}{\sigma^2}x-\frac{\mu^2}{2\sigma^2}\right\}\frac{\exp\left\{-\dfrac{x^2}{2\sigma^2}\right\}}{\sqrt{2\pi\sigma^2}}\notag\\
&= \exp\{\eta(\mu)T(x)-B(\mu)\}h(x).
\end{align}
Como $\eta(\mu)=\mu/\sigma^2$ é uma função crescente de $\mu$, então, pelo Teorema 1, o teste UMP é aquele que rejeita $\mathcal{H}_0$ se, e somente se, $\sum_{i=1}^n x_i>k^*$, em que a constante $k^*$ é dada pela solução de
\begin{equation}
\alpha=\P_{\mu_0}\left(\sum_{i=1}^n X_i>k^*\right)=1-\Phi\left(\frac{k^*-n\mu_0}{\sqrt{n}\sigma}\right).
\end{equation}
Portanto, $\dfrac{k^*-n\mu_0}{\sqrt{n}\sigma}=z_{1-\alpha}$, em que $z_{1-\alpha}$ corresponde ao $(1-\alpha)$-ésimo quantil da distribuição normal padrão.

Desse modo, o teste UMP rejeita $\mathcal{H}_0$ se $\sum_{i=1}^n x_i>n\mu_0+\sqrt{n}\sigma z_{1-\alpha}$ ou, de forma equivalente, rejeita $\mathcal{H}_0$ se $\bar{x}>\mu_0+\dfrac{\sigma}{\sqrt{n}}z_{1-\alpha}$.
\end{theorem}
% Correspondência: 10 - Testes de Hipóteses 2.pdf, p. 31-32/95

% ---------- Teorema/desenvolvimento: H0: mu <= mu0 vs H1: mu > mu0, sigma^2 desconhecido ----------
\begin{theorem}[$\mathcal{H}_0:\mu\leq\mu_0$ \emph{vs.} $\mathcal{H}_1:\mu>\mu_0$ -- $\sigma^2$ desconhecido]
Sejam $X_1,X_2,\ldots,X_n$ uma A.A. de uma população $X\sim\normdist{\mu,\sigma^2}$, com $\sigma^2$ desconhecido. Nesse caso, $\Theta=\{(\mu,\sigma^2)\in\mathbb{R}^2\mid-\infty<\mu<\infty;\sigma^2>0\}$ e $\Theta_0=\{(\mu,\sigma^2)\in\mathbb{R}^2\mid\mu\leq\mu_0;\sigma^2>0\}$.

No caso em que $(\mu,\sigma^2)\in\Theta$, temos que os EMV's de $\mu$ e $\sigma^2$ correspondem a $\hat{\mu}=\bar{X}$ e $\hat{\sigma}^2=\sum_{i=1}^n(X_i-\bar{X})^2/n$, respectivamente. Assim,
\begin{align}
\sup_{(\mu,\sigma^2)\in\Theta}L(\mu,\sigma^2\mid\v{x}) &= L(\hat{\mu},\hat{\sigma}^2\mid\v{x}) = \frac{1}{(2\pi\hat{\sigma}^2)^{n/2}}\exp\left[-\frac{1}{2\hat{\sigma}^2}\sum_{i=1}^n(x_i-\hat{\mu})^2\right]\notag\\
&= \left[\frac{n}{2\pi\sum_{i=1}^n(x_i-\bar{x})^2}\right]^{n/2}\exp\left(-\frac{n}{2}\right).
\end{align}
Quando $(\mu,\sigma^2)\in\Theta_0$, pode-se mostrar (mostre!) que os EMV's dos parâmetros serão
\begin{equation}
\tilde{\mu} = \begin{cases}
\bar{X}, & \text{se } \bar{X}\leq\mu_0;\\
\mu_0, & \text{se } \bar{X}>\mu_0
\end{cases}
\quad \text{e} \quad
\tilde{\sigma}^2 = \begin{cases}
\dfrac{1}{n}\sum_{i=1}^n\left(X_i-\bar{X}\right)^2, & \text{se } \bar{X}\leq\mu_0;\\
\dfrac{1}{n}\sum_{i=1}^n(X_i-\mu_0)^2, & \text{se } \bar{X}>\mu_0.
\end{cases}
\end{equation}
Desse modo,
\begin{align}
\sup_{(\mu,\sigma^2)\in\Theta_0}L(\mu,\sigma^2\mid\v{x}) &= L(\tilde{\mu},\tilde{\sigma}^2\mid\v{x}) = \frac{1}{(2\pi\tilde{\sigma}^2)^{n/2}}\exp\left[-\frac{1}{2\tilde{\sigma}^2}\sum_{i=1}^n(x_i-\tilde{\mu})^2\right]\notag\\
&= \begin{cases}
\left[\dfrac{n}{2\pi\sum_{i=1}^n(x_i-\bar{x})^2}\right]^{n/2}\exp\left(-\dfrac{n}{2}\right), & \text{se } \bar{x}\leq\mu_0;\\[2ex]
\left[\dfrac{n}{2\pi\sum_{i=1}^n(x_i-\mu_0)^2}\right]^{n/2}\exp\left(-\dfrac{n}{2}\right), & \text{se } \bar{x}>\mu_0.
\end{cases}
\end{align}
A razão de verossimilhanças será
\begin{align}
\lambda(\v{x}) &= \frac{\sup_{(\mu,\sigma^2)\in\Theta_0}L(\mu,\sigma^2\mid\v{x})}{\sup_{(\mu,\sigma^2)\in\Theta}L(\mu,\sigma^2\mid\v{x})}\notag\\
&= \begin{cases}
1, & \text{se } \bar{x}\leq\mu_0;\\[1ex]
\left[\dfrac{\sum_{i=1}^n(x_i-\bar{x})^2}{\sum_{i=1}^n(x_i-\mu_0)^2}\right]^{n/2}, & \text{se } \bar{x}>\mu_0.
\end{cases}
\end{align}
Antes de continuarmos, observe que
\begin{align}
\frac{\sum_{i=1}^n(x_i-\bar{x})^2}{\sum_{i=1}^n(x_i-\mu_0)^2} &= \frac{\sum_{i=1}^n(x_i-\bar{x})^2}{\sum_{i=1}^n(x_i-\bar{x}+\bar{x}-\mu_0)^2} = \frac{\sum_{i=1}^n(x_i-\bar{x})^2}{\sum_{i=1}^n(x_i-\bar{x})^2+n(\bar{x}-\mu_0)^2}\notag\\
&= \frac{1}{1+n\dfrac{(\bar{x}-\mu_0)^2}{\sum_{i=1}^n(x_i-\bar{x})^2}} = \frac{1}{1+n\dfrac{(\bar{x}-\mu_0)^2}{(n-1)s^2}} = \frac{1}{1+\dfrac{1}{n-1}\dfrac{(\bar{x}-\mu_0)^2}{s^2/n}}\notag\\
&= \left[1+\frac{1}{n-1}\frac{(\bar{x}-\mu_0)^2}{s^2/n}\right]^{-1} = \left(1+\frac{t^2}{n-1}\right)^{-1},
\end{align}
em que $t=\dfrac{\bar{x}-\mu_0}{s/\sqrt{n}}$. Logo,
\begin{equation}
\lambda(\v{x}) = \begin{cases}
1, & \text{se } \bar{x}\leq\mu_0;\\[1ex]
\left(1+\dfrac{t^2}{n-1}\right)^{-n/2}, & \text{se } \bar{x}>\mu_0.
\end{cases}
\end{equation}
Desse modo, para $0<\lambda_0<1$, rejeitamos $\mathcal{H}_0$ se, e somente se,
\begin{equation}
\lambda(\v{x})\leq\lambda_0 \iff t\geq k.
\end{equation}
Denotemos $T=\dfrac{\bar{X}-\mu_0}{S/\sqrt{n}}$. Sabemos que, se $\mu=\mu_0$, então $T\sim\tdist{n-1}$. Portanto, o valor de $k$ é tal que
\begin{equation}
\alpha=\P_{\mu_0}(T\geq k)=\P_{\mu_0}\left(T\geq t_{n-1,1-\alpha}\right),
\end{equation}
em que $t_{n-1,1-\alpha}$ corresponde ao $(1-\alpha)$-ésimo quantil de uma $\tdist{n-1}$.

O teste rejeita $\mathcal{H}_0$ se, e somente se,
\begin{equation}
\frac{\bar{x}-\mu_0}{s/\sqrt{n}}\geq t_{n-1,1-\alpha},
\end{equation}
ou, de forma equivalente,
\begin{equation}
\bar{x}\geq\mu_0+\frac{s}{\sqrt{n}}t_{n-1,1-\alpha}.
\end{equation}
\end{theorem}
% Correspondência: 10 - Testes de Hipóteses 2.pdf, p. 33-39/95

% ---------- Teorema/desenvolvimento: H0: mu = mu0 vs H1: mu != mu0, sigma^2 desconhecido ----------
\begin{theorem}[$\mathcal{H}_0:\mu=\mu_0$ \emph{vs.} $\mathcal{H}_1:\mu\neq\mu_0$ -- $\sigma^2$ desconhecido]
Sejam $X_1,X_2,\ldots,X_n$ uma A.A. de uma população $X\sim\normdist{\mu,\sigma^2}$, com $\sigma^2$ desconhecido. No caso em que $\sigma^2$ for desconhecido, o espaço paramétrico irrestrito e o espaço paramétrico sob $\mathcal{H}_0$ serão $\Theta=\{(\mu,\sigma^2)\in\mathbb{R}^2\mid-\infty<\mu<\infty;\sigma^2>0\}$ e $\Theta_0=\{(\mu,\sigma^2)\in\mathbb{R}^2\mid\mu=\mu_0;\sigma^2>0\}$.

Já vimos que
\begin{equation}
\sup_{(\mu,\sigma^2)\in\Theta}L(\mu,\sigma^2\mid\v{x}) = \left[\frac{n}{2\pi\sum_{i=1}^n(x_i-\bar{x})^2}\right]^{n/2}\exp\left(-\frac{n}{2}\right).
\end{equation}
Quando $(\mu,\sigma^2)\in\Theta_0$, não é difícil verificar (verifique!) que os EMV's dos parâmetros são, respectivamente,
\begin{equation}
\tilde{\mu}=\mu_0 \quad \text{e} \quad \tilde{\sigma}^2=\frac{1}{n}\sum_{i=1}^n(X_i-\mu_0)^2.
\end{equation}
Logo,
\begin{align}
\sup_{(\mu,\sigma^2)\in\Theta_0}L(\mu,\sigma^2\mid\v{x}) &= L(\tilde{\mu},\tilde{\sigma}^2\mid\v{x}) = \frac{1}{(2\pi\tilde{\sigma}^2)^{n/2}}\exp\left[-\frac{1}{2\tilde{\sigma}^2}\sum_{i=1}^n(x_i-\tilde{\mu})^2\right]\notag\\
&= \left[\frac{n}{2\pi\sum_{i=1}^n(x_i-\mu_0)^2}\right]^{n/2}\exp\left(-\frac{n}{2}\right).
\end{align}
De forma análoga ao que vimos anteriormente,
\begin{align}
\lambda(\v{x}) &= \frac{\sup_{(\mu,\sigma^2)\in\Theta_0}L(\mu,\sigma^2\mid\v{x})}{\sup_{(\mu,\sigma^2)\in\Theta}L(\mu,\sigma^2\mid\v{x})} = \left[\frac{\sum_{i=1}^n(x_i-\bar{x})^2}{\sum_{i=1}^n(x_i-\mu_0)^2}\right]^{n/2}\notag\\
&= \left(1+\frac{t^2}{n-1}\right)^{-n/2},
\end{align}
em que $t=\dfrac{\bar{x}-\mu_0}{s/\sqrt{n}}$.

Assim, para $0<\lambda_0<1$,
\begin{align}
\left(1+\frac{t^2}{n-1}\right)^{-n/2}\leq\lambda_0 &\iff\notag\\
|t|\geq\sqrt{(n-1)\left(\lambda_0^{-2/n}-1\right)}\equiv k &\iff\notag\\
t\leq-k \quad \text{ou} \quad t\geq k.
\end{align}
Denotemos $T=\dfrac{\bar{X}-\mu_0}{S/\sqrt{n}}$. Sabemos que, sob $\mathcal{H}_0$, $T\sim\tdist{n-1}$.

O valor de $k$ será solução de
\begin{align}
\alpha &= \P_{\mu_0}\left([T\leq-k]\text{ ou }[T\geq k]\right) = \P_{\mu_0}(T\leq-k)+\P_{\mu_0}(T\geq k)\notag\\
&= 2\P_{\mu_0}(T\geq k).
\end{align}
Portanto, $k=t_{n-1,1-\alpha/2}$.

Desse modo, o teste da RVG rejeita $\mathcal{H}_0$ se, e somente se,
\begin{equation}
t\leq-t_{n-1,1-\alpha/2} \quad \text{ou} \quad t\geq t_{n-1,1-\alpha/2},
\end{equation}
ou seja,
\begin{equation}
\bar{x}\leq\mu_0-t_{n-1,1-\alpha/2}\frac{s}{\sqrt{n}} \quad \text{ou} \quad \bar{x}\geq\mu_0+t_{n-1,1-\alpha/2}\frac{s}{\sqrt{n}}.
\end{equation}
Vale destacar que esse teste não é UMP. Por outro lado, pode-se mostrar que é um teste não viesado uniformemente mais poderoso.
\end{theorem}
% Correspondência: 10 - Testes de Hipóteses 2.pdf, p. 41-47/95

% ---------- Testes para a variância: introdução ----------
Assim como feito para testes para a média, vamos abordar testes unilaterais do tipo $\mathcal{H}_0:\sigma^2\leq\sigma_0^2$ \emph{vs.} $\mathcal{H}_1:\sigma^2>\sigma_0^2$, bem como testes bilaterais.

Há dois casos a serem considerados:
\begin{outline}
\1 $\mu$ conhecido;
\1 $\mu$ desconhecido.
\end{outline}
% Correspondência: 10 - Testes de Hipóteses 2.pdf, p. 48/95

% ---------- Teorema/desenvolvimento: H0: sigma^2 <= sigma0^2 vs H1: sigma^2 > sigma0^2, mu conhecido ----------
\begin{theorem}[$\mathcal{H}_0:\sigma^2\leq\sigma_0^2$ \emph{vs.} $\mathcal{H}_1:\sigma^2>\sigma_0^2$ -- $\mu$ conhecido]
Sejam $X_1,X_2,\ldots,X_n$ uma A.A. de uma população $X\sim\normdist{\mu,\sigma^2}$, com $\mu$ conhecido. Nesse caso, o espaço paramétrico será um intervalo (conjunto dos reais positivos). Além disso,
\begin{align}
f(x\mid\mu) &= \frac{1}{\sqrt{2\pi\sigma^2}}\exp\left\{-\frac{1}{2\sigma^2}(x-\mu)^2\right\}\notag\\
&= \exp\left\{-\frac{1}{2\sigma^2}(x-\mu)^2-\frac{1}{2}\log\sigma^2\right\}\frac{1}{\sqrt{2\pi}}\notag\\
&= \exp\{\eta(\sigma^2)T(x)-B(\sigma^2)\}h(x).
\end{align}
Seja $R=\sum_{i=1}^n\left(\dfrac{X_i-\mu}{\sigma_0}\right)^2$. Já sabemos que, se a verdadeira variância for $\sigma_0^2$, então $R\sim\chidist{n}$.

Observe que $\eta(\sigma^2)=-1/(2\sigma^2)$ é uma função crescente de $\sigma^2$. Pelo Teorema 1, o teste UMP rejeita $\mathcal{H}_0$ se, e somente se, $\sum_{i=1}^n(x_i-\mu)^2>k^*$, em que $k^*$ é solução de
\begin{align}
\alpha=\P_{\sigma_0^2}\left(\sum_{i=1}^n(X_i-\mu)^2>k^*\right) &= \P_{\sigma_0^2}\left(R>k^*/\sigma_0^2\right) = \P_{\sigma_0^2}\left(R>\chi^2_{n,1-\alpha}\right).
\end{align}
Portanto, $\dfrac{k^*}{\sigma_0^2}=\chi^2_{n,1-\alpha}$.

Assim, o teste UMP rejeita $\mathcal{H}_0$ se, e somente se, $\sum_{i=1}^n(x_i-\mu)^2>\sigma_0^2\chi^2_{n,1-\alpha}$.
\end{theorem}
% Correspondência: 10 - Testes de Hipóteses 2.pdf, p. 49-50/95

% ---------- Teorema/desenvolvimento: H0: sigma^2 = sigma0^2 vs H1: sigma^2 != sigma0^2, mu desconhecido ----------
\begin{theorem}[$\mathcal{H}_0:\sigma^2=\sigma_0^2$ \emph{vs.} $\mathcal{H}_1:\sigma^2\neq\sigma_0^2$ -- $\mu$ desconhecido]
Sejam $X_1,X_2,\ldots,X_n$ uma A.A. de uma população $X\sim\normdist{\mu,\sigma^2}$, com $\mu$ desconhecido. Nesse caso, $\Theta=\{(\mu,\sigma^2)\in\mathbb{R}^2\mid-\infty<\mu<\infty;\sigma^2>0\}$ e $\Theta_0=\{(\mu,\sigma^2)\in\mathbb{R}^2\mid-\infty<\mu<\infty;\sigma^2=\sigma_0^2\}$.

Já vimos que
\begin{equation}
\sup_{(\mu,\sigma^2)\in\Theta}L(\mu,\sigma^2\mid\v{x}) = \left(\frac{1}{2\pi\hat{\sigma}^2}\right)^{n/2}\exp\left(-\frac{n}{2}\right).
\end{equation}
Quando $(\mu,\sigma^2)\in\Theta_0$, não é difícil verificar (verifique!) que os EMV's dos parâmetros são, respectivamente,
\begin{equation}
\tilde{\mu}=\bar{X} \quad \text{e} \quad \tilde{\sigma}^2=\sigma_0^2.
\end{equation}
Logo,
\begin{align}
\sup_{(\mu,\sigma^2)\in\Theta_0}L(\mu,\sigma^2\mid\v{x}) &= L(\tilde{\mu},\tilde{\sigma}^2\mid\v{x}) = \frac{1}{(2\pi\tilde{\sigma}^2)^{n/2}}\exp\left[-\frac{1}{2\tilde{\sigma}^2}\sum_{i=1}^n(x_i-\tilde{\mu})^2\right]\notag\\
&= \frac{1}{(2\pi\sigma_0^2)^{n/2}}\exp\left[-\frac{1}{2\sigma_0^2}\sum_{i=1}^n(x_i-\bar{x})^2\right] = \frac{1}{(2\pi\sigma_0^2)^{n/2}}\exp\left(-\frac{n\hat{\sigma}^2}{2\sigma_0^2}\right).
\end{align}
Assim,
\begin{align}
\lambda(\v{x}) &= \frac{\sup_{(\mu,\sigma^2)\in\Theta_0}L(\mu,\sigma^2\mid\v{x})}{\sup_{(\mu,\sigma^2)\in\Theta}L(\mu,\sigma^2\mid\v{x})} = \left(\frac{\hat{\sigma}^2}{\sigma_0^2}\right)^{n/2}\exp\left(-\frac{n}{2}\frac{\hat{\sigma}^2}{\sigma_0^2}+\frac{n}{2}\right)\notag\\
&\equiv y^{n/2}\exp\left[-\frac{n}{2}(y-1)\right],
\end{align}
em que $y=\dfrac{\hat{\sigma}^2}{\sigma_0^2}$.

Assim, para $0<\lambda_0<1$,
\begin{align}
y^{n/2}\exp\left[-\frac{n}{2}(y-1)\right]\leq\lambda_0 &\iff\notag\\
y\leq k_1 \quad \text{ou} \quad y\geq k_2, & 
\end{align}
em que $y=\dfrac{\hat{\sigma}^2}{\sigma_0^2}$.

Seja $V=\dfrac{\sum_{i=1}^n\left(X_i-\bar{X}\right)^2}{\sigma_0^2}$. Nesse caso, sob $\mathcal{H}_0$, $V\sim\chidist{n-1}$. Além disso, note que $\dfrac{\hat{\sigma}^2}{\sigma_0^2}=\dfrac{V}{n}$.

Assim, os valores de $k_1$ e $k_2$ são soluções de
\begin{align}
\alpha &= \P_{\sigma_0^2}\left(\left[\frac{\hat{\sigma}^2}{\sigma_0^2}\leq k_1\right]\text{ ou }\left[\frac{\hat{\sigma}^2}{\sigma_0^2}\geq k_2\right]\right) = \P_{\sigma_0^2}\left(\left[\frac{V}{n}\leq k_1\right]\text{ ou }\left[\frac{V}{n}\geq k_2\right]\right)\notag\\
&= \P_{\sigma_0^2}(V\leq nk_1)+\P_{\mu_0}(V\geq nk_2)\notag\\
&= \P_{\sigma_0^2}\left(V\leq\chi^2_{n-1,\alpha/2}\right)+\P_{\mu_0}\left(V\geq\chi^2_{n-1,1-\alpha/2}\right) \quad \text{(uma possibilidade).}
\end{align}
Nesse caso mais simples (que não é o melhor!), temos
\begin{equation}
k_1=\frac{\chi^2_{n-1,\alpha/2}}{n} \quad \text{e} \quad k_2=\frac{\chi^2_{n-1,1-\alpha/2}}{n}.
\end{equation}
Desse modo, o teste rejeita $\mathcal{H}_0$ se, e somente se,
\begin{equation}
\frac{\hat{\sigma}^2}{\sigma_0^2}\leq\frac{\chi^2_{n-1,\alpha/2}}{n} \quad \text{ou} \quad \frac{\hat{\sigma}^2}{\sigma_0^2}\geq\frac{\chi^2_{n-1,1-\alpha/2}}{n},
\end{equation}
em que $\chi^2_{r,\beta}$ corresponde ao $\beta$-ésimo quantil de uma distribuição $\chidist{r}$.
\end{theorem}
% Correspondência: 10 - Testes de Hipóteses 2.pdf, p. 53-59/95

% ---------- Texto: distribuição assintótica do TRVG ----------
À medida que o número de parâmetros do modelo aumenta, torna-se mais complicada a tarefa de estabelecer testes da RVG.

Por vezes, essa tarefa não é tão simples nem mesmo em situações uniparamétricas.

Uma forma de contornar esse problema é através de aproximação.
% Correspondência: 10 - Testes de Hipóteses 2.pdf, p. 61/95

% ---------- Teorema 3 ----------
\begin{theorem}
Sejam $X_1,X_2,\ldots,X_n$ uma A.A. de uma população com f.d.p. (ou f.p.) $f(x\mid\theta)$, em que $\theta$ pode ser escalar ou vetor. Sob certas condições de regularidade, se $\theta\in\Theta_0$, então
\begin{equation}
-2\log\Lambda_n \toD \chidist{r},
\end{equation}
quando $n\to\infty$. O número de graus de liberdade $r$ corresponde ao número de restrições em $\Theta_0$.
\end{theorem}
% Correspondência: 10 - Testes de Hipóteses 2.pdf, p. 62/95

% ---------- Observação 8 ----------
\begin{note}
	\leavevmode
	\begin{outline}[enumerate]
		\1 O número de graus de liberdade $r$ também pode ser interpretado como a dimensão de $\Theta$ menos a dimensão de $\Theta_0$.
	\end{outline}
\end{note}
% Correspondência: 10 - Testes de Hipóteses 2.pdf, p. 62/95

% ---------- Texto: princípio da RVG no caso assintótico ----------
Lembrando que $\Lambda_n$ é uma v.a. que assume valores do tipo
\begin{equation}
\lambda_n=\frac{\sup_{\theta\in\Theta_0}L(\theta\mid\v{x})}{\sup_{\theta\in\Theta}L(\theta\mid\v{x})},
\end{equation}
que corresponde à RVG para uma amostra observada $\v{x}=(x_1,x_2,\ldots,x_n)$.

O princípio da RVG sugere que rejeitemos $\mathcal{H}_0$ para valores pequenos de $\lambda_n$. Portanto, no caso assintótico, o teste deve rejeitar a hipótese nula para valores grandes de $-2\log\lambda_n$.
% Correspondência: 10 - Testes de Hipóteses 2.pdf, p. 63/95

Levando em consideração a distribuição assintótica quando $\mathcal{H}_0$ for verdadeira, um teste com tamanho aproximadamente $\alpha$ rejeita $\mathcal{H}_0$ se, e somente se,
\begin{equation}
-2\log\lambda_n\geq\chi^2_{r,1-\alpha},
\end{equation}
em que $\chi^2_{r,1-\alpha}$ corresponde ao $(1-\alpha)$-ésimo quantil de uma distribuição $\chidist{r}$.
% Correspondência: 10 - Testes de Hipóteses 2.pdf, p. 64/95

